\section*{Chapter \ref{ch:rc:liquidity-and-funding-risk-bank-treasury} --- Liquidity and Funding Risk; Bank Treasury}

\begin{solution}{exo:rc:liquidity-and-funding-risk-bank-treasury:1}
Level~2A after the haircut is $0.85\times40 = 34$, capped at $\tfrac23\times30 = 20$ (Level~2 at most 40\% of the
stock); HQLA $= 30+20 = 50$.
\end{solution}

\begin{solution}{exo:rc:liquidity-and-funding-risk-bank-treasury:2}
$-6.2\times3.5\% = -21.7\%$ to first order; the exact fall is smaller (19.2\% for the chapter's MBS) because
of positive convexity, and because the duration itself falls as yields rise.
\end{solution}

\begin{solution}{exo:rc:liquidity-and-funding-risk-bank-treasury:3}
$300\times e^{-5/4} = 86$ basis points.
\end{solution}

\begin{solution}{exo:rc:liquidity-and-funding-risk-bank-treasury:4}
Its assets are longer than its liabilities (long fixed-rate securities against deposits slotted over
five years at most): it is long duration, like a receiver of fixed in a swap. It gains when rates fall
and loses when they rise.
\end{solution}

\begin{solution}{exo:rc:liquidity-and-funding-risk-bank-treasury:5}
Uninsured depositors lose money if the bank fails, so any doubt about its solvency gives each of them a
reason to leave first; when they are few, large and connected, they leave together. Their run-off is
fast and correlated, which is a liquidity risk to the bank whatever its capital.
\end{solution}

\begin{solution}{exo:rc:liquidity-and-funding-risk-bank-treasury:6}
The business is charged the term funding rate of its loans (4.84\% for five years in the example), so a
long loan earns only its credit and service margin; the term and liquidity premium goes to the treasury,
which funds and hedges it. Without FTP the business would book the curve's slope as profit.
\end{solution}

\begin{solution}{exo:rc:liquidity-and-funding-risk-bank-treasury:7}
$+0.86$ a year with a beta of 0.35, $-1.87$ with 0.80.
\end{solution}

\begin{solution}{exo:rc:liquidity-and-funding-risk-bank-treasury:8}
Holding to maturity avoids realising the loss only if the bank never needs to sell or pledge the
securities; their economic value has fallen all the same, the bank earns below-market yields on them
for years, and a deposit outflow can force a sale. Government guarantees remove credit risk, not
interest rate or liquidity risk.
\end{solution}

\begin{solution}{pb:rc:liquidity-and-funding-risk-bank-treasury:1}
\textbf{1.} USD~17.2 billion, 108\% of equity.
\textbf{2.} USD~3.9 billion.
\textbf{3.} $16-17.2-3.9 = -5.1$: negative.
\textbf{4.} A 200 basis point rise already cut EVE by 6.95, 43\% of Tier~1, far above the 15\% outlier
threshold: the duration mismatch was visible before any loss.
\textbf{5.} Pay-fixed swaps of matching duration: EVE immune to parallel moves, but paying fixed at
market rates on the notional reduced the carry of the securities to about the floating rate.
\textbf{6.} 93\%: the uninsured deposits carry a 40\% run-off, and the MBS count only after the haircut
and within the Level~2 cap.
\textbf{7.} USD~36.1 billion, 20.9\% of deposits.
\textbf{8.} SVB lost over 40 billion dollars on 9 March, about a quarter of the roughly 165 billion
implied by its management's figures (140 billion being 85\%); the stylised bank could not meet that.
\textbf{9.} Repo of the Treasuries and MBS, the discount window against collateral at market value less
haircuts, Federal Home Loan Bank advances; each takes time and marks the securities.
\textbf{10.} It lent against securities at par, not market value, for up to a year: the bank could borrow
on its MBS without realising the 19\% loss.
\textbf{11.} Retail deposits are insured, many and dispersed, and move slowly; uninsured corporate deposits
are large, concentrated and informed, and move together.
\textbf{12.} NII falls as deposit costs rise while fixed assets do not reprice; EVE falls further because
the deposits are worth less as a cheap funding source.
\textbf{13.} That the bank had to sell at a loss and needed capital: a signal of weakness that invited
the run.
\textbf{14.} Because the run concerns liquidity at market value: the depositors can be paid only from
assets whose market value is below their book value, and each depositor fears being last.
\textbf{15.} None of the standard ratios: concentration and connectedness of depositors need their own
metrics (share of uninsured, largest depositors, sector concentration).
\textbf{16.} Hedge the duration of the securities, keep them shorter, price deposits with a realistic
beta, and hold liquidity for the uninsured base.
\textbf{17.} At least in supervisory measures and \omterm{def:rc:stress-testing-and-scenarios:stress}{stress tests}: capital that ignores them overstates
the loss-absorbing capacity in a run.
\textbf{18.} Against history by cohort and rate cycle, against peers, with stress scenarios of rapid
outflows, and by reviewing the concentration of the base.
\textbf{19.} Named result: \emph{March 2023}: the stylised bank's HTM unrealised loss is 108\% of its
equity, and its liquid assets run out after a one-day outflow of 20.9\% of deposits (USD~36.1 billion).
\textbf{20.} Both: market risk made the bank insolvent at market value, and liquidity risk turned that
into failure in a day.
\end{solution}

\begin{solution}{iq:rc:liquidity-and-funding-risk-bank-treasury:1}
LCR: HQLA against 30 days of stressed net outflows, a short-term survival buffer. NSFR: stable funding
against the stable funding needs of assets over a year, a limit on maturity transformation.
\iqlookfor{the two horizons and what each constrains.}
\end{solution}

\begin{solution}{iq:rc:liquidity-and-funding-risk-bank-treasury:2}
An internal price for funds by tenor (base curve plus liquidity premium) charged to asset businesses
and credited to liability businesses; it measures margins net of funding, centralises rate and
liquidity risk in treasury, and prices liquidity use.
\iqlookfor{the mechanism and the incentives.}
\end{solution}

\begin{solution}{iq:rc:liquidity-and-funding-risk-bank-treasury:3}
EVE values all future cash flows; NII measures the next year. A bank with long fixed assets and
sticky deposits loses EVE when rates rise (the assets fall) but gains NII for a while (floating
assets reprice, deposits do not).
\iqlookfor{horizon difference and an example.}
\end{solution}

\begin{solution}{iq:rc:liquidity-and-funding-risk-bank-treasury:4}
Balances split into core and volatile, with decay rates estimated from history by segment; deposit
rates as a function of market rates (beta, lags, floors); stress overlays for outflows; validation by
cycle, including rate rises.
\iqlookfor{balances, rates and validation.}
\end{solution}

\begin{solution}{iq:rc:liquidity-and-funding-risk-bank-treasury:5}
Long-duration securities funded by concentrated, uninsured deposits, unhedged rate risk, reliance on
held-to-maturity accounting, weak liquidity stress testing, and a public announcement that triggered a
coordinated run.
\iqlookfor{duration, funding concentration and governance.}
\end{solution}

\begin{solution}{iq:rc:liquidity-and-funding-risk-bank-treasury:6}
Pay-fixed swaps against the securities or receive-floating on the deposits' modelled duration; this
reduces NII when curves are upward sloping, and needs hedge accounting to avoid volatility in reported
earnings.
\iqlookfor{instrument, carry cost and accounting.}
\end{solution}
